\chapter{Hash-Based Signatures}
\label{ch:hash-based}

\abstract*{Hash-based signatures achieve quantum resistance using only the assumption that hash functions are one-way---the most minimal and conservative foundation in all of post-quantum cryptography. This chapter traces the evolution from Lamport's one-time signatures through Merkle trees to the stateless SPHINCS+ scheme standardized by NIST as SLH-DSA, developing each construction with full algorithmic detail.}

% =============================================================
% SOURCE: notes.tex §3.4.3 (lines 3945–4236)
% REUSE: ~20%. Currently one subsection. Mostly new chapter.
% =============================================================

\section{One-Time Signatures}
\label{sec:ots}

% TODO: Adapt from notes.tex §3.4.3 (Lamport)
% - Lamport signatures: construction, correctness, security proof
% - Key and signature sizes: the efficiency problem
% - Winternitz OTS (WOTS+): trading size for computation
% - WOTS+ parameters: Winternitz parameter w, chains
% - Cite: Lamport 1979

\section{Merkle's Tree-Based Signatures}
\label{sec:merkle}

% TODO: Expand from notes.tex §3.4.3
% - From one-time to many-time: the Merkle tree construction
% - Authentication paths
% - Key generation, signing, verification algorithms
% - The state management requirement
% - Cite: Merkle 1979

\section{Stateful Schemes: XMSS and LMS}
\label{sec:stateful}

% TODO: New section
% - XMSS: eXtended Merkle Signature Scheme
% - Multi-tree XMSS (XMSS^MT)
% - LMS: Leighton-Micali Signature scheme
% - State management: the critical operational challenge
% - Standards: RFC 8391 (XMSS), RFC 8554 (LMS), NIST SP 800-208
% - When stateful schemes are appropriate: firmware signing, CA roots

\section{FORS: Few-Time Signatures from Random Subsets}
\label{sec:fors}

% TODO: New section
% - The FORS construction: key idea
% - Security analysis: few-time signature bounds
% - Role in SPHINCS+ as the bottom layer

\section{SLH-DSA (SPHINCS+)}
\label{sec:slh-dsa}

% TODO: New section (the main event of this chapter)
% - The hypertree construction: layers of XMSS trees
% - FORS at the bottom, WOTS+ at internal nodes
% - Removing state: the randomized message hash trick
% - Parameter sets: SLH-DSA-128s/f, -192s/f, -256s/f
% - "s" (small) vs "f" (fast) parameter tradeoffs
% - Security proof: reduction to hash function properties
% - Cite: FIPS 205, Bernstein et al. (SPHINCS+ paper)

\section{Security Analysis}
\label{sec:hash-security}

% TODO: New section
% - Minimal assumptions: one-wayness, second-preimage resistance
% - Multi-target attacks and parameter adjustments
% - Grover's impact: doubling hash output lengths
% - Comparison with lattice-based assumptions: strength of foundation
% - The quantum random oracle model for hash-based schemes

%% Exercises
\begin{prob}
\label{prob:ch10-lamport}
Implement a Lamport one-time signature scheme using SHA-256. Sign a 256-bit message, verify the signature, then demonstrate why the scheme is insecure if the same key is used to sign two different messages.
\end{prob}

\begin{prob}
\label{prob:ch10-merkle}
Construct a Merkle tree of height $h = 3$ (supporting 8 one-time signatures). Draw the tree, show the authentication path for the 5th leaf, and explain the verification procedure.
\end{prob}

\begin{prob}
\label{prob:ch10-wots}
In WOTS+ with Winternitz parameter $w = 16$ and $n = 256$ (security parameter in bits), compute: (a)~the number of chains, (b)~the signature size in bytes, (c)~the average number of hash evaluations for signing and verification. Compare with Lamport for the same security level.
\end{prob}
